% =============================================================================
%                              MODULE OVERVIEW
% =============================================================================
\section*{Module Overview}
\addcontentsline{toc}{section}{Module Overview}

% -----------------------------------------------------------------------------
%                              INTRODUCTION
% -----------------------------------------------------------------------------
\subsection*{Introduction}

Ethical issues relating to the use and deployment of current and future AI technologies are becoming more and more of a concern, and are regularly discussed and argued about in the media. The urgency with which these issues are discussed is in large part due to recent dramatic advances in machine learning (ML) approaches to AI, such as large language models, and the envisaged extent to which AI machines may be autonomous and act against our interests.

These ML advances are in turn driven by the increasing availability of big data and its use by machine learning algorithms, which are also exploited by social media companies to profile and predict our behaviours. We therefore see concerns about the deployment of autonomous cars and weapons, the possibility of super-intelligent machines and how their actions might inadvertently cause harm, and the problems of data bias and the use of data in targeted advertising.

This module reviews basic concepts in philosophy that are particularly relevant to how we should think about and address ethical concerns related to the AI algorithms of today and of the future.

\separator

% -----------------------------------------------------------------------------
%                                 AIMS
% -----------------------------------------------------------------------------
\subsection*{Aims}

On successful completion of this module, you will be able to demonstrate knowledge of:
\begin{itemize}
  \item Basic philosophical work on the nature of \textbf{intelligence}, \textbf{consciousness}, \textbf{communication}, \textbf{society}, and \textbf{ethics}
  \item How this work relates to and impacts on developments in Artificial Intelligence
  \item The \textbf{ethical implications} of these developments
\end{itemize}

\separator

% -----------------------------------------------------------------------------
%                           LEARNING OUTCOMES
% -----------------------------------------------------------------------------
\subsection*{Learning Outcomes}

By the end of this module, you will be able to:

\begin{itemize}
  \item[\textbf{PEA1:}] Discuss basic philosophical work on the nature of intelligence, consciousness, communication, society, and ethics
  \item[\textbf{PEA2:}] Explain how this work relates to and impacts on developments in Artificial Intelligence
  \item[\textbf{PEA3:}] Apply critical reasoning skills to challenging materials
  \item[\textbf{PEA4:}] Engage in an informative way with issues of public concern regarding the dangers of Artificial Intelligence
\end{itemize}

\separator

% -----------------------------------------------------------------------------
%                            COURSE STRUCTURE
% -----------------------------------------------------------------------------
\subsection*{Course Structure}

\begin{center}
\renewcommand{\arraystretch}{1.3}
\begin{tabular}{@{} c l p{7cm} @{}}
\toprule
\textbf{Week} & \textbf{Topic} & \textbf{Key Themes} \\
\midrule
1 & Introduction & Module overview, philosophical foundations \\
2 & Intelligence & What is intelligence? Turing Test, Chinese Room \\
3 & Consciousness & Mind-body problem, qualia, machine consciousness \\
4 & Reasoning \& Communication & Logic, language, meaning, understanding \\
5 & Ethics \& Morality & Moral theories, ethical frameworks for AI \\
6 & Algorithms & Algorithmic decision-making, bias, fairness \\
7 & AI in Medicine \& War & Applications, risks, ethical considerations \\
8 & Superintelligence & Value loading, alignment, existential risk \\
9 & AI \& Human Society & Social impact, future of work, governance \\
10 & Deceptive AI & (Not examined) \\
11 & Review/Revision & Exam preparation \\
\bottomrule
\end{tabular}
\end{center}

\separator

% -----------------------------------------------------------------------------
%                         KEY PHILOSOPHERS & CONCEPTS
% -----------------------------------------------------------------------------
\subsection*{Key Philosophers and Concepts to Know}

This module draws on work from many philosophers. Key figures likely to appear include:

\begin{itemize}
  \item \textbf{Alan Turing} --- Turing Test, computability, machine intelligence
  \item \textbf{John Searle} --- Chinese Room argument, intentionality
  \item \textbf{Thomas Nagel} --- ``What is it like to be a bat?'', consciousness
  \item \textbf{David Chalmers} --- Hard problem of consciousness
  \item \textbf{Immanuel Kant} --- Deontological ethics, categorical imperative
  \item \textbf{John Stuart Mill} --- Utilitarianism, consequentialism
  \item \textbf{Nick Bostrom} --- Superintelligence, existential risk
\end{itemize}

\textit{Note: This list will be expanded as we cover each week's content.}

